\subsubsection{Fermion Hierarchies from Fricke Modular Symmetry --- arXiv:2610.29754}

\textbf{Authors:} Takafumi Kai, Justin Kaidi, Katsumi Kina, Hajime Otsuka

\paragraph{主な主張}
Frickeモジュラー群 $\Gamma_0^+(3)$ の固定点近傍で単一のmodulusの変位からクォーク質量のべき指数 $(7,3,0)$・$(5,4,2)$ とCKM混合の $(1,2,3)$ を実現し、実数のYukawa係数と複素VEVからCKMのCP位相を再現した\cite{Kai:2026Fricke}。

\paragraph{新規性}
Fricke拡張によって内部固定点の残存対称性を $\mathbb{Z}_3$ から $\mathbb{Z}_6$ へ拡大し、Fricke parityが展開の先頭次数を追加で3次抑制することを、カスプで正則なモジュラー形式の完全な空間について示した。

\paragraph{位置付け}
残存モジュラー対称性によるフレーバー階層生成をFricke群へ拡張した具体例であり、必要なmodulusの真空を実現する安定化機構は今後の課題として残る。
